\chapter{اللقاء الرابع عشر: تابع تفسير سورة البقرة (٣)}

\section{مقدمة: اغتنام الأوقات في طلب العلم}
السلام عليكم ورحمة الله وبركاته. الحمد لله رب العالمين، والصلاة والسلام على نبينا محمد وعلى آله وصحبه أجمعين. هذا هو الدرس الثاني عشر (الرابع عشر بترقيم دورتنا) من اجتماعنا على مدارسة تفسير \rawi{ابن جرير الطبري} رحمه الله.

أعتذر عن التأخر في الأيام الماضية لانشغالي، وهذا هو السبب في أنني أحياناً أطيل الدروس أو أعطي دروساً متتالية؛ لأن العاقل هو الذي يغتنم فراغه قبل شغله. فكلما وجدت فرصة لإنجاز شيء من البرنامج بادرت به ولا أؤخره، خشية أن يطرأ طارئ يشغل الإنسان. ونسأل الله تبارك وتعالى أن يرزقنا الهدى والتقى ويعيننا على إنجاز ما نطلب من الخير.

\separator

\section{تأويل (أفتطمعون أن يؤمنوا لكم)}
\isnad{قال الإمام الطبري رحمه الله:} «القول في تأويل قوله تعالى: \begin{ayah}أَفَتَطْمَعُونَ أَن يُؤْمِنُوا لَكُمْ\end{ayah}... أي أفترجون يا معشر المؤمنين بمحمد ﷺ... أن يؤمن لكم يهود بني إسرائيل... أن يصدقوكم بما جاءكم به نبيكم».

\begin{fawaid}[تعقيب على تفسير الإيمان بالتصديق]
نلاحظ أن \rawi{الطبري} رحمه الله يفسر الإيمان هنا بالتصديق. والذي أراه والله أعلم أن الإيمان أعلى من مجرد التصديق، بل فيه معنى الائتمان من جهة (في الخبر)، ومعنى الاتباع والتسليم من جهة أخرى (في الأمر). فلما قال اليهود لموسى: \textit{لن نؤمن لك حتى نرى الله جهرة}، ليس المراد التصديق المجرد، بل لن نتبعك ولن نبقى على الإيمان بك.
\end{fawaid}

\isnad{وأما قوله تعالى:} \begin{ayah}وَقَدْ كَانَ فَرِيقٌ مِّنْهُمْ\end{ayah} \isnad{قال الطبري:} «إنما جعل الله الذين كانوا على عهد موسى ومن بعدهم من بني إسرائيل من اليهود (منهم)... لأنهم كانوا آباءهم وأسلافهم».

\separator

\section{تأويل (يسمعون كلام الله ثم يحرفونه) وعقوبات الإعراض}
اختلف المفسرون في الذين عنى الله أنهم \begin{ayah}يَسْمَعُونَ كَلَامَ اللَّهِ\end{ayah}، فرأى \rawi{مجاهد} و\rawi{السدي} أن المراد (يسمعون التوراة ثم يحرفونها). ورأى آخرون (كـ \rawi{الربيع بن أنس} و\rawi{ابن إسحاق}) أن المراد الذين باشروا سماع كلام الله (حين سألوا موسى رؤية الله فأخذتهم الصاعقة). 

ورجح \rawi{الطبري} القول الثاني، محتجاً بأنه لو كان المراد التوراة لكان كل اليهود يسمعونها، فما وجه تخصيص (فريق منهم)؟

\begin{fawaid}[عقوبة رد الآيات البينات]
في هذه القصة عبرة؛ وهي أن الآيات البينات إذا عُرضت على العبد فكذبها وعاندها، فإن من عقوبات ذلك أن يُحال بينه وبين الإيمان ويُختم على قلبه، كما قال تعالى: \textit{ونقلب أفئدتهم وأبصارهم كما لم يؤمنوا به أول مرة ونذرهم في طغيانهم يعمهون}. فكل من جاءته السنن فصد عنها، يوشك أن يُصرف عن الإيمان.
\end{fawaid}

\separator

\section{تأويل (وإذا خلا بعضهم إلى بعض قالوا أتحدثونهم بما فتح الله عليكم)}
\isnad{قال الطبري رحمه الله:} «وإذا خلا بعض هؤلاء اليهود إلى بعض... قالوا أتحدثونهم بما فتح الله عليكم؟».
واختلف المفسرون في معنى الفتح هنا: فمنهم من قال هو (ما أمركم الله به)، ومنهم من قال (بما في كتابكم من نعت محمد ﷺ)، ومنهم من قال هو (بما فتح الله عليكم من العذاب).

\begin{fawaid}[مهارة التلخيص وبناء الخرائط الذهنية]
من الفواصل بين القارئ العادي وطالب العلم الجاد هي مهارة "التلخيص". حين تسرد الأقوال، يجب أن تبني خريطة ذهنية بخلاصات الأقوال، كأن تستخرج أسهماً للقول الأول (نعت محمد) والقول الثاني (العذاب)، ثم تلخص حجة الطبري في الترجيح بأسلوبك أو باختصار من عبارته، فهذا يدل على الفهم ويُثبّت العلم.
\end{fawaid}

ورجح \rawi{الطبري} القول بأن الفتح هنا هو: «بما حكم الله به عليكم وقضاه فيكم... ومن حكمه ما أخذ به ميثاقهم من الإيمان بمحمد ﷺ». لأن أصل الفتح في كلام العرب: النصر والقضاء والحكم.

\isnad{وأما قوله:} \begin{ayah}أَوَلَا يَعْلَمُونَ أَنَّ اللَّهَ يَعْلَمُ مَا يُسِرُّونَ وَمَا يُعْلِنُونَ\end{ayah} فهو توبيخ لليهود الذين ظنوا أن إسرارهم بكتمان الحق وتلاومهم عليه يخفى على الله.

\begin{fawaid}[العلم بالله مفتاح الهداية والطمأنينة]
القرآن يعلمنا عن الله، وكلما علمت عن الله أسمائه وأفعاله ازددت هداية وطمأنينة. فهذه الآية وإن كانت تهديداً للكفار، إلا أنها طمأنينة للمؤمن أن ربه محيط بالكافرين وليس بغافل عما يعمل الظالمون.
\end{fawaid}

\separator

\section{تأويل (ومنهم أميون) والاستدراك على معناها}
\isnad{قال أبو جعفر:} «يعني بالأميين الذين لا يكتبون ولا يقرأون... ومنه قول النبي ﷺ: إنا أمة لا نكتب ولا نحسب... وروي عن ابن عباس قول خلاف هذا القول... قال: الأميون قوم لم يصدقوا رسولاً... وهذا التأويل خلاف ما يُعرف من كلام العرب المستفيض».

\begin{fawaid}[مناقشة في معنى الأمي]
يظهر لي -والله أعلم- أن تفسير "الأمي" هنا بمجرد الذي لا يقرأ ولا يكتب يحتاج إلى مراجعة. فالأمي في القرآن استُعمل كثيراً في مقابلة (أهل الكتاب)، فالأميون هم العرب الذين لم ينزل عليهم كتاب من قبل، كما كان اليهود يتكبرون ويقولون (ليس علينا في الأميين سبيل)، وكقوله تعالى (هو الذي بعث في الأميين رسولاً). فقوله (ومنهم أميون) كأنه يشير إلى قوم دخلوا في دين اليهود من العرب وليس عندهم علم الكتاب كالأحبار.
\end{fawaid}

\subsection{تأويل (لا يعلمون الكتاب إلا أماني) والاستثناء المنقطع}
رجح \rawi{الطبري} تفسير \rawi{مجاهد} للأماني بأنها: الأكاذيب والتقوّل على الله (من التخرص والتمني باختلاق الكذب).

\begin{fawaid}[الاستثناء المنقطع]
إذا قيل: كيف يُستثنى الكذب من الكتاب في قوله (لا يعلمون الكتاب إلا أماني) والكذب ليس من القرآن؟ الجواب: هذا يسميه أهل العربية "الاستثناء المنقطع"، وهو أن يأتي الاستثناء بـ (إلا) بمعنى (لكن)، ليخرج ما بعدها عن معنى وصنف ما قبلها. أي: لا يعلمون الكتاب، لكنهم يختلقون الأكاذيب.
\end{fawaid}

\separator

\section{تأويل (فويل للذين يكتبون الكتاب بأيديهم) وإبطال دعوى الزيادة}
\isnad{قال أبو جعفر:} «يعني بذلك جل ثناؤه أحبار اليهود الذين حرفوا كتاب الله... وكتبوا كتاباً على ما تأولوه... والويل العذاب في الوادي السائل من صديد أهل النار».

\isnad{وطرح الطبري إشكالاً:} فما وجه قوله \textit{بأيديهم} وهل يكتبون بغير اليد؟
\isnad{وأجاب:} «أعلم ربنا جل ثناؤه قوله ذلك... أن أحبار اليهود تلي الكذب والفرية على الله بأيديهم على علم منهم... فنفى الله بقوله (يكتبون الكتاب بأيديهم) أن يكون ولي كتابة ذلك بعض جهالهم بأمر علمائهم... بل باشروا ذلك بأيديهم».

\begin{fawaid}[بطلان الزيادة في القرآن بلا معنى]
يرد بعض اللغويين والمتكلمين كلمات كـ (بأيديهم) أو (يقولون بألسنتهم) أو (يطير بجناحيه) ويزعمون أنها زائدة لا تفيد معنى جديداً. وهذا باطل؛ فالقرآن ليس فيه حرف زيادة بلا دلالة، بل كل لفظ له فائدة تقريرية أو نفي لاحتمال مجازي (كالنفي هنا أن يكونوا أمروا غيرهم بالكتابة بل باشروها بأيديهم). وقد صنف ابن تيمية رسالة جليلة في نفي هذا الادعاء.
\end{fawaid}

\separator

\section{تأويل (بلى من كسب سيئة وأحاطت به خطيئته) والرد على الخوارج}
\isnad{قال أبو جعفر:} «تكذيب من الله جل ثناؤه القائلين من اليهود: لن تمسنا النار إلا أياماً معدودة... وإنما قلنا إن السيئة التي ذكرها الله عز وجل أن من كسبها وأحاطت به خطيئته فهو من أهل النار المخلدين فيها، إنما عنى جل ذكره بها بعض السيئات دون بعض... والخلود في النار لأهل الكفر بالله دون أهل الإيمان به لتظاهر الأخبار عن رسول الله ﷺ أن أهل الإيمان لا يخلدون فيها».

\begin{fawaid}[تخصيص اللفظ العام بالقرينة]
هذا موضع دقيق في منهج الطبري، وهو أنه خصص اللفظ العام (السيئة) ببعض أفرادها وهو (الشرك). وحجته في ذلك السياق والقرينة؛ وهي أن الله رتب على هذه السيئة الخلود الأبدي في النار (فأولئك أصحاب النار هم فيها خالدون)، والخلود لا يكون إلا للمشركين. وهذا من أحسن الأمثلة على تخصيص العام بحجة.
\end{fawaid}

\begin{fawaid}[الرد على المعتزلة والخوارج في الشفاعة والخلود]
ضمن الطبري كتابه هذا الرد على المبتدعة كالخوارج والمعتزلة الذين يستدلون بمثل هذه الآيات القائلة بالخلود للمسيء، أو بنفي الشفاعة (ولا يقبل منها شفاعة)، على تخليد مرتكب الكبيرة الموحد في النار! وهذا استدلال باطل يخالف الكتاب والسنة المتواترة في أن الشفاعة تنال أهل الكبائر من الموحدين، وأن الخلود منفي عنهم.
\end{fawaid}

\separator

\section{خاتمة: نصائح منهجية في طلب العلم}
يا شباب، طلب العلم ليس مجرد حشو للمعلومات في العقل، بل هو افتقار إلى الله في طلب الهدى، واستعانة به، وتحلٍّ بالصبر والحلم والأناة وحسن الخلق.

احذروا من سرعة الإعجاب والاندفاع، فمن يسرع الإعجاب بأحد يسرع الهدم له لاحقاً إن رأى منه زلة. واختاروا بعناية من تدرسون معهم ومن تستمعون إليهم؛ فمجالس العلم ليست كمنصات وسائل التواصل التي يتكلم فيها من لا علم له ولا خلق له. واجعلوا ميزانكم دائماً كتاب الله وسنة نبيه، ومسلك الصحابة الكرام. 

والسلام عليكم ورحمة الله وبركاته.